\section[Realización física discreta sobre rutas TPK]{Realización física
discreta: lecturas positiva y orientada sobre rutas TPK, dominio común y
localización reversible}
\label{sec:int68-realizacion-fisica-discreta}

La forma espectral, la realización de Hilbert, la sección dimensional y el
transporte de Cartan no son decoraciones añadidas a una salida clásica. Son
cuatro representaciones coordinadas de una misma ruta TPK enriquecida. Para
expresar esa unidad causal debe construirse primero su dominio común; no
basta escribir cuatro flechas con el mismo símbolo de partida.

\subsection{Dominio común como producto fibrado}

Sea \(\mathsf P_{\rm TPK}^{\rm enr}\) la categoría de estados TPK preparados
y rutas registradas compatibles con sus extremos, frontera y restricciones.
Sean
\[
 \mathsf D_{\rm sp}:=\mathsf{SurvPath}_{\rm HMT},\qquad
 \mathsf D_Q:=\operatorname{Dom}(\mathcal Q),\qquad
 \mathsf D_D:=\operatorname{Dom}(\mathfrak D_{\rm HMT}),\qquad
 \mathsf D_C:=\operatorname{Dom}(\rho_C^{\rm disc}),
\]
con sus funtores de olvido fieles
\[
 U_i:\mathsf D_i\longrightarrow\mathsf P_{\rm TPK}^{\rm enr}
 \qquad(i\in\{{\rm sp},Q,D,C\}).
\]
Cada \(U_i\) olvida únicamente la realización y conserva el estado y la ruta
TPK subyacentes. Definimos
\begin{equation}
 \mathsf P_{\rm CT}^{+}
 :=
 \mathsf D_{\rm sp}
 \times_{\mathsf P_{\rm TPK}^{\rm enr}}\mathsf D_Q
 \times_{\mathsf P_{\rm TPK}^{\rm enr}}\mathsf D_D
 \times_{\mathsf P_{\rm TPK}^{\rm enr}}\mathsf D_C.
 \label{eq:int68-pct-pullback}
\end{equation}
Un objeto es, por tanto, un estado preparado provisto de las cuatro
realizaciones sobre el mismo antecedente; una flecha es una ruta registrada
cuya imagen subyacente coincide en los cuatro factores. Si las
presentaciones se identifican sólo por isomorfismo, se usa el producto
fibrado bicategórico correspondiente; en la presentación fijada del registro
se adopta el producto fibrado estricto de
\eqref{eq:int68-pct-pullback}.

\begin{proposition}[Existencia de un objeto en el dominio común]
\label{prop:int68-pct-no-vacio}
La categoría \(\mathsf P_{\rm CT}^{+}\) no es vacía. No es necesario
particularizar la ruta electrónica central para probarlo: cualquier sección
coinductiva certificada, equipada con su flecha identidad y sus datos neutros, admite
levantamientos compatibles a los cuatro dominios de
\eqref{eq:int68-pct-pullback}.
\end{proposition}

\begin{proof}
El \cref{thm:generacion-monodromica-conjunta-rev7} proporciona, por ejemplo,
una sección superviviente preparada \(x_\pi\). Su identidad
\(\operatorname{id}_{x_\pi}\) es simultáneamente: una ruta superviviente;
un morfismo de \(\mathsf P_{\rm TPK}\), que \(\mathcal Q\) envía a la
identidad acotada de \(\mathcal Q(x_\pi)\); la ruta vacía con firma nula,
que el \cref{thm:funtor-dimensional-minimo} envía a
\((0,0,0,1)\); y la identidad del grupoide, que el
\cref{thm:realizacion-discreta-cartan-rev6} envía al elemento neutro
\((I,0)\) de Poincaré. Estos cuatro levantamientos olvidan al mismo objeto
\(x_\pi\) y al mismo morfismo \(\operatorname{id}_{x_\pi}\). Su cuádrupla
define, por la propiedad del producto fibrado, un objeto de
\(\mathsf P_{\rm CT}^{+}\).
\end{proof}

\begin{definition}[Adaptadores al dominio común]
\label{def:int68-pct-positive}
Los cuatro adaptadores son las proyecciones canónicas
\[
 \begin{aligned}
 F_{\rm sp}:\mathsf P_{\rm CT}^{+}&\longrightarrow
     \mathsf{SurvPath}_{\rm HMT},\\
 F_Q:\mathsf P_{\rm CT}^{+}&\longrightarrow\operatorname{Dom}(\mathcal Q),\\
 F_D:\mathsf P_{\rm CT}^{+}&\longrightarrow
     \operatorname{Dom}(\mathfrak D_{\rm HMT}),\\
 F_C:\mathsf P_{\rm CT}^{+}&\longrightarrow
     \operatorname{Dom}(\rho_C^{\rm disc}).
 \end{aligned}
\]
En particular, \(F_{\rm sp}\) conserva la familia compatible de
restricciones finitas de la sección y de la ruta; no reconstruye esa familia
a partir de su espectro.
\end{definition}

\begin{lemma}[Funtorialidad de los adaptadores]
\label{lem:int68-adaptadores}
Los cuatro adaptadores preservan identidades y composición.
\end{lemma}

\begin{proof}
La identidad de un objeto del producto fibrado es la cuádrupla de identidades
sobre el mismo estado TPK; cada proyección devuelve la identidad del factor.
La composición está definida componente a componente entre cuádruplas que
proyectan sobre la misma concatenación de rutas. Proyectar antes o después de
componer da, por definición, el mismo resultado.
\end{proof}

\subsection{\(4\) componentes y producto positivo}

Los teoremas de realización espectral, hilbertiana, dimensional y de Cartan
proporcionan
\[
 \mathfrak S:\mathsf{SurvPath}_{\rm HMT}
      \longrightarrow\mathbf{SpecRoute},
 \qquad
 \mathcal Q:\operatorname{Dom}(\mathcal Q)
      \longrightarrow\mathsf{Hilb}_{\mathbb R,J},
\]
\[
 \mathfrak D_{\rm HMT}:\operatorname{Dom}(\mathfrak D_{\rm HMT})
      \longrightarrow\mathbf B\mathfrak M_D^+,
 \qquad
 \rho_C^{\rm disc}:\operatorname{Dom}(\rho_C^{\rm disc})
      \longrightarrow\mathbf B G_C,
\]
donde
\[
 G_C:=\operatorname{Spin}^{+}(1,3)\ltimes\mathbb R^{1,3}
\]
y
\begin{equation}
 \mathfrak M_D^+
 :=\mathbb N\times\mathbb R_{\geq0}^{2}\times\mathbb R_{>0},
 \qquad
 (n,a,b,\lambda)\odot(n',a',b',\lambda')
 =(n+n',a+a',b+b',\lambda\lambda').
 \label{eq:int68-monoide-dimensional-positivo}
\end{equation}
Aquí \(\mathbf BM\) es la categoría de un objeto asociada al monoide
\(M\). El primer funtor es el del
\cref{thm:functor-realizacion-espectral-numero}; el segundo es la
realización de Hilbert definida en el capítulo de postulados; el tercero es
el \cref{thm:funtor-dimensional-minimo}; y el cuarto es el
\cref{thm:realizacion-discreta-cartan-rev6}.

Definimos las restricciones comunes
\[
 \mathfrak S_+:=\mathfrak S\circ F_{\rm sp},\qquad
 \mathcal Q_+:=\mathcal Q\circ F_Q,\qquad
 \mathfrak D_+:=\mathfrak D_{\rm HMT}\circ F_D,\qquad
 \rho_{C,+}^{\rm disc}:=\rho_C^{\rm disc}\circ F_C.
\]

\begin{theorem}[Realización discreta positiva sobre el dominio común]
\label{thm:int68-producto-positivo}
La aplicación
\begin{equation}
 \mathfrak R_{\rm MD}^{\rm disc,+}
 :=(\mathfrak S_+,\mathcal Q_+,\mathfrak D_+,
       \rho_{C,+}^{\rm disc})
 \label{eq:int68-producto-positivo}
\end{equation}
es un funtor
\[
 \mathsf P_{\rm CT}^{+}\longrightarrow
 \mathbf{SpecRoute}\times\mathsf{Hilb}_{\mathbb R,J}
 \times\mathbf B\mathfrak M_D^+\times\mathbf BG_C.
\]
No se afirma que toda ruta TPK pertenezca a \(\mathsf P_{\rm CT}^{+}\).
\end{theorem}

\begin{proof}
Por el \cref{lem:int68-adaptadores}, cada proyección \(F_i\) es un funtor;
por los teoremas citados, también lo son
\(\mathfrak S,\mathcal Q,\mathfrak D_{\rm HMT}\) y
\(\rho_C^{\rm disc}\) en sus dominios declarados. Sus composiciones
preservan, pues, identidades y concatenación. La composición en el codominio
producto se calcula componente a componente, lo que prueba
\eqref{eq:int68-producto-positivo}.
\end{proof}

La componente de Hilbert tiene por blanco general
\(\mathsf{Hilb}_{\mathbb R,J}\), cuyos morfismos son operadores reales
acotados que entrelazan las estructuras complejas. Sea
\(\mathsf P_{\rm CT}^{u}\subseteq\mathsf P_{\rm CT}^{+}\) la subcategoría
de las flechas \(\gamma:x\to y\) para las cuales
\begin{equation}
 \mathcal Q_+(\gamma)^*\mathcal Q_+(\gamma)=I_{\mathcal Q_+(x)},
 \qquad
 \mathcal Q_+(\gamma)\mathcal Q_+(\gamma)^*=I_{\mathcal Q_+(y)}.
 \label{eq:int68-unitariedad-doble}
\end{equation}
Las identidades y los productos de operadores unitarios vuelven a ser
unitarios; por ello esta condición define una subcategoría y sólo en ella se
obtiene
\begin{equation}
 \mathcal U_{\rm reg}:=
 \mathcal Q_+\big|_{\mathsf P_{\rm CT}^{u}}:
 \mathsf P_{\rm CT}^{u}\longrightarrow
 \mathsf{UHilb}_{\mathbb R,J}.
 \label{eq:int68-restriccion-unitaria}
\end{equation}
La invertibilidad unilateral o la conservación de norma en un subconjunto no
sustituyen las dos identidades de \eqref{eq:int68-unitariedad-doble}.

La notación de siete entradas se descompone en las cuatro componentes
tipadas: \(Q\) es \(\mathcal Q_+\);
\(U\) es la restricción \eqref{eq:int68-restriccion-unitaria}; \(A\) es la
coordenada de acción de \(\mathfrak D_+\); \(C\) es la conjugación declarada
en el subdominio donde entrelaza la componente hilbertiana y la orientación;
\(\mathcal E\) y \(\mathcal B\) son lectoras constitutivas subordinadas a
los datos dimensionales; y \(N_{\rm sp}\) es una propiedad cuadrática de la
realización partida. Ninguna de las cuatro últimas se añade al producto como
un quinto funtor sin dominio y codominio.

\subsection{Lecturas positiva y orientada}

Sea \(e\mapsto e^\vee\) la reversión de aristas del grafo dirigido
subyacente. En \(\mathsf P_{\rm CT}^{+}\), \(e^\vee\) es otra arista y no
se declara todavía inversa categorial de \(e\). Fijado un peso
\(w(e)=w(e^\vee)>0\), definimos sobre una ruta
\(\gamma=e_1\cdots e_r\)
\[
 V_+(\gamma):=\sum_{j=1}^r w(e_j),
 \qquad
 V_{\rm or}(\gamma):=
 \sum_{j=1}^r\varepsilon(e_j)w(\underline e_j),
\]
donde \(\underline e_j\) es la arista no orientada y
\(\varepsilon(e^\vee)=-\varepsilon(e)\). La primera lectora conserva
variación total positiva; la segunda, transporte neto orientado. Ambas son
aditivas por concatenación y reciben la misma ruta TPK.

\begin{proposition}[No factorización del lector positivo por el orientado]
\label{prop:int68-no-factorizacion}
Supóngase que la identidad y el lazo
\(ee^\vee:=e^\vee\circ e\) pertenecen al dominio construido. No existe una
función \(H\) sobre \(\operatorname{im}V_{\rm or}\) tal que
\(V_+=H\circ V_{\rm or}\).
\end{proposition}

\begin{proof}
Se tiene
\[
 V_{\rm or}(\operatorname{id})=0
 =V_{\rm or}(ee^\vee),
 \qquad
 V_+(\operatorname{id})=0,
 \qquad
 V_+(ee^\vee)=2w(e)>0.
\]
Una misma entrada \(0\) obligaría a \(H(0)\) a tomar dos valores distintos.
\end{proof}

Este resultado es una no factorización concreta en el dominio construido; no
es un teorema universal sobre cualquier entrelazador enriquecido posible.
En particular, no autoriza a borrar una lectora. La unidad está en el
antecedente TPK común y no en una compresión horizontal entre sus salidas.

\subsection{Localización reversible y realización orientada}

Sea \(\mathsf P_{{\rm CT},{\rm rev}}^+\) la subcategoría generada por las
aristas para las cuales se han comprobado las cuatro condiciones siguientes:
\begin{enumerate}
 \item \(\mathfrak S_+(e)\) es un isomorfismo de
       \(\mathbf{SpecRoute}\) y
       \(\mathfrak S_+(e^\vee)=\mathfrak S_+(e)^{-1}\);
 \item \(\mathcal Q_+(e)\) satisface
       \eqref{eq:int68-unitariedad-doble} y
       \(\mathcal Q_+(e^\vee)=\mathcal Q_+(e)^*\);
 \item el dato dimensional orientado procede de una cocadena aditiva con
       \(d_{\rm or}(e^\vee)=-d_{\rm or}(e)\) y toma valores en
       \(\operatorname{Gr}(\mathfrak M_D^+)\);
 \item los cociclos de Cartan son normalizados, aditivos y antisimétricos
       bajo reversión.
\end{enumerate}
Formamos el grupoide
\[
 \Pi_\vee(\mathsf P_{{\rm CT},{\rm rev}}^+)
\]
obtenido de la completación en grupoide imponiendo además
\([e^\vee]=[e]^{-1}\). La componente dimensional orientada toma valores en
\begin{equation}
 \operatorname{Gr}(\mathfrak M_D^+)
 \cong\mathbb Z\times\mathbb R^2\times\mathbb R_{>0},
 \label{eq:int68-grothendieck-dimensional}
\end{equation}
con inversa
\((n,a,b,\lambda)^{-1}=(-n,-a,-b,\lambda^{-1})\). Este lector firmado no es
la magnitud positiva \(V_+\) rebautizada: sus valores sobre las aristas
reversas son distintos.

Para la componente Cartan, sean \(c_g,c_s\) cociclos como en el punto 4,
sea \(R_4\in\operatorname{Spin}^{+}(1,3)\), sea \(R_4u=u\) y sea
\(\ell_0>0\). Entonces
\[
 \rho_C^{\rm disc}(\gamma)
 =\bigl(R_4^{c_g(\gamma)},\ell_0c_s(\gamma)u\bigr).
\]
La ley
\((R,a)(S,b)=(RS,a+Rb)\), la igualdad \(R_4u=u\) y la aditividad de los
cociclos dan
\[
 \rho_C^{\rm disc}(\gamma_2\gamma_1)
 =\rho_C^{\rm disc}(\gamma_2)\rho_C^{\rm disc}(\gamma_1),
 \qquad
 \rho_C^{\rm disc}(\gamma^\vee)
 =\rho_C^{\rm disc}(\gamma)^{-1}.
\]

\begin{theorem}[Realización discreta orientada]
\label{thm:int68-realizacion-orientada}
Bajo las cuatro condiciones anteriores, las componentes descienden de manera
única a un funtor
\[
 \mathfrak R_{\rm MD}^{\rm disc,or}:
 \Pi_\vee(\mathsf P_{{\rm CT},{\rm rev}}^+)
 \longrightarrow
 \operatorname{Core}(\mathbf{SpecRoute})
 \times\mathsf{UHilb}_{\mathbb R,J}
 \times\mathbf B\operatorname{Gr}(\mathfrak M_D^+)
 \times\mathbf BG_C.
\]
En particular,
\[
 \mathfrak R_{\rm MD}^{\rm disc,or}(\gamma^{-1})
 =\mathfrak R_{\rm MD}^{\rm disc,or}(\gamma)^{-1}.
\]
\end{theorem}

\begin{proof}
Las condiciones 1 y 2 hacen invertibles las imágenes espectral e
hilbertiana; la condición 3 toma valores en un grupo; y la componente Cartan
ya toma valores en \(G_C\). Las cuatro respetan la relación generadora
\([e^\vee]=[e]^{-1}\). Por la propiedad universal de la localización en
grupoide, cada una desciende de manera única; su producto es el funtor
indicado. La identidad de inversión se verifica componente a componente.
\end{proof}

La lectora positiva \(\mathfrak R_{\rm MD}^{\rm disc,+}\) y la orientada
\(\mathfrak R_{\rm MD}^{\rm disc,or}\) son hermanas: ninguna se define como
compresión de la otra. El paso posterior de un valor orientado a su módulo no
es, en general, monoidal y coincide con una lectura positiva sólo en un cono
sin cancelación. Del mismo modo, este par no sustituye la doble proyección
fundacional: la evaluación arquimediana y la rama
Paley--Witt--Golay--Leech--\(V^\natural\)--Monstruo--Moonshine permanecen
como lectoras hermanas del mismo estado enriquecido. Las cuatro componentes
construidas en esta sección no sustituyen ninguna de esas dos proyecciones.

\subsection{Entrelazamiento de Cartan y límite suave}

Sea
\[
 \mathcal A_{\ell_0}
 =\begin{pmatrix}\omega&\ell_0^{-1}e\\0&0\end{pmatrix}
\]
una conexión de Cartan declarada. La ecuación
\begin{equation}
 \operatorname{Hol}_{\mathcal A_{\ell_0}}
   \bigl(\mathfrak R_C(\gamma)\bigr)
 =\rho_C\bigl(\operatorname{Hol}_{\rm TPK}(\gamma)\bigr)
 \label{eq:int68-cuadrado-cartan}
\end{equation}
es exacta en el modelo discreto cuando ambas holonomías se definen por el
mismo producto ordenado de aristas.

\begin{criterion}[Promoción a una realización suave]
\label{crit:int68-cartan-suave}
La ecuación \eqref{eq:int68-cuadrado-cartan} sólo define una prolongación
suave si se prueban, además, independencia de subdivisión, convergencia del
producto ordenado y recuperación local de una cotétrada y una conexión con
los dominios y regularidad declarados.
\end{criterion}

El término «físico» designa aquí que las cuatro componentes publican
espectro, realización de estado, dimensiones y transporte de Cartan sin
perder su antecedente TPK\@. No afirma unicidad, equivalencia de categorías,
cobertura de toda ruta, carta SI automática ni teoría de campos completa.

\begin{criterion}[Obstrucciones a las realizaciones positiva, unitaria y orientada]
\label{crit:int68-obstrucciones-realizacion}
La realización positiva no está definida si falta alguno de los cuatro
adaptadores, si sus imágenes no poseen la misma ruta TPK subyacente o si una
componente no preserva identidad y composición. La restricción unitaria no
existe si falta cualquiera de las dos igualdades de
\eqref{eq:int68-unitariedad-doble}. La realización orientada no desciende al
grupoide si una imagen espectral no es isomorfa, si una relación localizada
cambia de imagen, si una magnitud positiva se trata como inversa, si una
reversión no cambia el signo de los datos orientados o si el cociclo de
Cartan deja de ser normalizado, aditivo o antisimétrico. No existe
prolongación suave compatible si el resultado depende de la subdivisión o si
el límite no recupera cotétrada y conexión. Cualquier uso de una coordenada
externa para seleccionar la ruta viola el dominio construido.
\end{criterion}

\paragraph{Dominio, codominio y estatuto de las ecuaciones de confluencia HMT–MD}
El producto positivo queda demostrado sobre el producto fibrado común y la
realización orientada, sobre la localización reversible bajo las cuatro
condiciones explícitas. Extender el dominio a otras rutas exige construir sus
cuatro levantamientos sobre un mismo antecedente. La prolongación suave exige,
además, las propiedades de subdivisión, convergencia y recuperación local del
\cref{crit:int68-cartan-suave}; una verificación finita no sustituye esas
propiedades.
